\section{Background: the predicates in question}
\label{sec:background}

For authorization there is no intrinsic oracle: vulnerable and correct
executions both answer with success. Metamorphic testing replaces the absent
oracle with a relation over multiple executions---and the relation is only as
sound as the assumptions it needs, which for authorization are exactly the facts
the oracle was supposed to supply. \mbox{MST-wi}~\cite{Chaleshtari2022Metamorphic},
building on Mai et al.~\cite{Mai2019Metamorphic}, crawls a per-user model of
reachable URLs and contents, then checks MRs against that recorded model. In the
artifact we inspected, the predicates guarding an authorization MR read either
the recorded model directly (\texttt{cannotReachThroughGUI},
\texttt{userCanRetrieveContent}) or static subject properties; we call a
predicate \emph{GUI-derived} when its value is a function of the crawl.

Two mechanical counts frame what follows. Over the \ResCatalogMrTotal{} catalog
MRs, \ResCatalogGuiDerived{} (\ResCatalogGuiDerivedPct{}\%) gate on a
GUI-derived or identity predicate, and only \ResCatalogUnion{}
(\ResCatalogUnionPct{}\%) are authorization-relevant by name or action---an
upper bound, reported as a caveat. And a full-text search for predicates reading
the \emph{state of an authorization artifact} (ownership, permission record,
share) finds \ResCatalogAuthzStateHits{} occurrences. The delegation to
reachability is therefore unavoidable, and the reading moment enters the
specification: if an entitlement changes after the crawl, the MR evaluates
reachability against a world that no longer exists. This is a statement about
what the relations \emph{mean}; our experiments are built so that the crawl-time
and live readings can be told apart.
